% GENERATED FILE. Edit canonical lessons, then run npm run generate:books.
% canonical-source-hash: fnv1a64:68bd391dcae0f5b2
% canonical-lessons: HI-C79-dopahar-raat-ka-khana

\chapter{Lunch and Dinner}
\label{ch:hi-dopahar-raat-ka-khana}
\hlchaptermodality{87}

\begin{chapteropening}
\textbf{By the end of this chapter:} \emph{I can name all three meals of the day, and say which one is a borrowed word and which two are phrases Hindi builds from a time plus khana.}

Say all three meals in order, morning to night, and name the odd one out: nashta is borrowed whole, while dopahar ka khana and rat ka khana are built from a time the learner already has plus khana.
\end{chapteropening}

\section[dopahar kā khānā / rāt kā khānā]{\hi{दोपहर} \hi{का} \hi{खाना} / \hi{रात} \hi{का} \hi{खाना} — the two meals you already know the words for}

\label{lesson:HI-C79-dopahar-raat-ka-khana}



\begin{quote}
You have \textbf{\hi{खाना}} (\emph{khānā}), \textquotedblleft{}food\textquotedblright{}. You have \textbf{\hi{दोपहर}} (\emph{dopahar}) and \textbf{\hi{रात}} (\emph{rāt}). Before reading on, guess what \textbf{\hi{दोपहर} \hi{का} \hi{खाना}} is.
\end{quote}

\subsection*{You'll want to know: \hi{दोपहर} \hi{का} \hi{खाना} and \hi{रात} \hi{का} \hi{खाना}}
\textbf{\hi{दोपहर} \hi{का} \hi{खाना}} (\emph{dopahar kā khānā}) — \textquotedblleft{}lunch.\textquotedblright{} \textbf{\hi{रात} \hi{का} \hi{खाना}} (\emph{rāt kā khānā}) — \textquotedblleft{}dinner.\textquotedblright{}

You guessed them, and that is the point of putting them here rather than in a word list. Neither is a new word. Each is \textbf{a time you already have, plus \hi{का}, plus the food-word you already have} — \textquotedblleft{}the afternoon's food\textquotedblright{}, \textquotedblleft{}the night's food.\textquotedblright{}

\begin{cousinweb}
The transparency runs deeper than the compound.

\textbf{\hi{दोपहर}} is itself two pieces you can see: \textbf{\hi{दो}} (\textquotedblleft{}two\textquotedblright{}) + \textbf{\hi{पहर}}, a traditional watch of roughly three hours. Two watches into the day is noon — the word is a clock reading that turned into a time of day.

\textbf{\hi{रात}} comes from Sanskrit \textbf{\hi{रात्रि}} (\emph{rātri}), the formal word that shows up in \textbf{\hi{शुभ} \hi{रात्रि}} (\textquotedblleft{}good night\textquotedblright{}). The everyday \textbf{\hi{रात}} is the same word worn down by use, which is why the greeting sounds more ceremonious than the noun.

So a three-word phrase for \textquotedblleft{}dinner\textquotedblright{} contains a number, an obsolete time unit, and a Sanskrit noun — and a learner who has met all three can read it cold.
\end{cousinweb}

\begin{grammarlens}[title={Grammar Lens: the pattern, not the pair}]
Set the day's three meals beside each other:

\noindent
\begin{tabularx}{\linewidth}{@{}>{\raggedright\arraybackslash}X>{\raggedright\arraybackslash}X>{\raggedright\arraybackslash}X@{}}
\toprule
\textbf{meal} & \textbf{Hindi} & \textbf{what it is built from} \\
\midrule
breakfast & \hi{नाश्ता} & its own word, borrowed whole \\
lunch & \hi{दोपहर} \hi{का} \hi{खाना} & a time + \hi{का} + \hi{खाना} \\
dinner & \hi{रात} \hi{का} \hi{खाना} & a time + \hi{का} + \hi{खाना} \\
\bottomrule
\end{tabularx}

\textbf{\hi{नाश्ता}} is the odd one out: a single borrowed noun with no parts to see. The other two are phrases Hindi builds on demand. That asymmetry is worth noticing because it keeps recurring — Hindi borrows a whole word where it has one, and builds a transparent phrase where it does not.
\end{grammarlens}

\subsection*{Your turn}
Say these aloud:

\begin{itemize}
\raggedright
  \item \textquotedblleft{}\hi{दोपहर} \hi{का} \hi{खाना}\textquotedblright{} — \emph{dopahar kā khānā}, lunch
  \item \textquotedblleft{}\hi{रात} \hi{का} \hi{खाना}\textquotedblright{} — \emph{rāt kā khānā}, dinner
  \item all three meals in order, morning to night (\hi{नाश्ता}, \hi{दोपहर} \hi{का} \hi{खाना}, \hi{रात} \hi{का} \hi{खाना})
\end{itemize}

\begin{tcolorbox}[breakable,colback=teal!4,colframe=teal!35!black,title={Before you move on}]
Which of the three meal names is a single borrowed word rather than a built phrase? (\textbf{\hi{नाश्ता}}.) What is \hi{दोपहर} made of, and what does it add up to? (\textbf{\hi{दो}} \textquotedblleft{}two\textquotedblright{} + \textbf{\hi{पहर}} a three-hour watch — two watches in, so noon.)
\end{tcolorbox}
